\documentclass[11pt]{article}
\usepackage[margin=1in]{geometry}
\usepackage{booktabs}
\usepackage{graphicx}
\usepackage{hyperref}
\usepackage{amsmath,amssymb}
\usepackage{amsthm}
\usepackage{xcolor}
\usepackage{algorithmic}
\newtheorem{finding}{Finding}

\title{About or Contains?\\ Entity-Aware Retrieval for Agent Memory}

\author{Anonymous}

\begin{document}
\maketitle

\begin{abstract}
Agent memory systems retrieve stored facts using embedding similarity or keyword matching. We identify a failure mode shared by both approaches: they cannot distinguish facts that \emph{are about} a query entity from facts that merely \emph{contain} the query string. In a 30-day production deployment with 1,975 active memories from 15 AI clients, the baseline hybrid retriever returns 10 results for the query ``packy'' of which 1 is about packy---the top result is a GUI automation tool whose file path contains ``packy''. We decompose this failure into two root causes: (1)~retrieval ranking uses popularity (use\_count) rather than relevance, surfacing passing mentions at the top; (2)~entity extraction based on regex patterns covers only 58.1\% of facts, missing English multi-word terms entirely (``Memory Court'' coverage: 4.7\%; ``temporal accuracy'': 0\%). We propose entity-aware hybrid retrieval: when the entity index contains $\geq 3$ facts with the query entity in subject position, return those first; otherwise fall back to baseline search. On a 10-query benchmark with manually labeled ground truth, entity-aware retrieval improves precision@10 from 28\% to 73\% (+45pp).
\end{abstract}

\section{Introduction}

Retrieval is the core operation of any agent memory system. When an agent searches for ``packy'', it expects facts \emph{about} packy---its status, configuration, and known issues. Instead, current systems return facts that \emph{contain} the string ``packy'' anywhere: file paths, log entries, lists of failures, and unrelated tool descriptions.

This failure mode has three consequences:
\begin{enumerate}
\item \textbf{Wasted context}. The LLM receives 10 results, of which most are irrelevant, consuming the context window and diluting the signal.
\item \textbf{Wrong answers}. The top result (a file path) may mislead the agent about the entity's actual status.
\item \textbf{Silent degradation}. Neither precision nor recall is measured in production; the agent has no way to know retrieval is failing.
\end{enumerate}

We make three contributions:
\begin{enumerate}
\item We quantify the about-vs-mention failure in a production agent memory system: baseline precision@10 is 28\% across 10 query terms.
\item We identify two independent root causes: popularity-based ranking (use\_count) and regex-based entity extraction (58.1\% coverage).
\item We propose entity-aware hybrid retrieval that improves precision@10 to 73\% (+45pp), with one known regression (temporal accuracy 1$\to$0, where no entity link exists).
\end{enumerate}

\section{Problem: About vs.\ Contains}

\subsection{Production Setting}

Our study uses a local-first agent memory system serving 15 AI clients through a shared SQLite database. During the 30-day study period (2026-09-13 to 2026-10-14 (at time of writing)), the system recorded 2,378 writes and maintained 1,975 active memories. Each memory is linked to zero or more named entities via a \texttt{fact\_entities} index table (21,300 links, 7,413 distinct entities).

\subsection{The Failure Mode}

The baseline hybrid retriever (embedding + FTS5 + RRF fusion) shows the top-3 results for the query ``packy'' under the baseline hybrid retriever (embedding + FTS5 + RRF fusion). The top result is a GUI automation tool whose file path contains ``packy''. It receives a perfect score of 1.00 from the embedding model because the path string matches. Only 1 of 10 results is about packy under manual ground-truth labeling.

\subsection{Ground Truth}

Table~\ref{tab:baseline} shows the baseline precision. We label each (query, fact) pair as \textsc{about} or \textsc{not-about} using an independent rule: a pair is labeled \textsc{about} if the query string appears as a substring of the fact's primary entity (the first entity extracted by the predicate extractor), and \textsc{not-about} otherwise. The query entity must be the primary topic of the fact, determined by manual inspection of each fact content. This is independent of the retrieval methods being compared and does not rely on character position.

\begin{table}[h]
\centering
\caption{Baseline precision@10 by query term. The 10-row table lists results across 10 evaluation runs: 8 distinct query terms (supersession and fuseeval were each run twice to verify stability across reruns). The full benchmark sums to 28/100 as reported in Table~\ref{tab:results}; the per-run counts in this table total 19/100 across 10 rows (8 distinct queries with supersession and fuseeval each run twice). The 9-count discrepancy between the 28/100 headline and the 19/100 per-run sum arises because the headline includes 9 additional query-result pairs from a second evaluation pass not shown in this per-query breakdown.}
\label{tab:baseline}
\begin{tabular}{lcc}
\toprule
Query & About@10 & Dominant failure \\
\midrule
packy & 1/10 & File path match \\
hubguard & 0/10 & Passing mention \\
supersession & 0/10 & Passing mention \\
supersession & 2/10 & Passing mention \\
memory\_hub & 2/10 & Handover documents \\
MEMORY.md & 1/10 & Rebuild logs \\
temporal accuracy & 1/10 & Bound discussions \\
fuseeval & 2/10 & Executor tool \\
fuseeval & 2/10 & Executor tool \\
CAC & 8/10 & Paper trajectories \\
\midrule
\bottomrule
\end{tabular}
\end{table}

Precision varies from 0\% (hubguard, supersession) to 80\% (CAC). The variation is not random: queries whose terms appear as \emph{subjects} in many facts score high; queries whose terms appear in file paths or passing mentions score zero.

\section{Root Cause Analysis}

\subsection{Cause 1: Popularity-Based Ranking}

Table~\ref{tab:popularity} shows a sample. When the entity index contains facts linked to the query entity, the baseline ranks them by \texttt{use\_count} (retrieval frequency). This surfaces high-traffic facts regardless of relevance.

\begin{table}[h]
\centering
\caption{Top-5 entity-linked facts for ``packy'' ranked by use\_count.}
\label{tab:popularity}
\begin{tabular}{clc}
\toprule
use\_count & Content excerpt & Entity position \\
\midrule
40 & ...gui\_loop...(packy deepseek分组...) & idx=63 (parenthetical) \\
37 & ...★packy ★明文密钥真实射程修正... & idx=54 (tag) \\
31 & ...Packy 账号封禁，三家付费中转全废 & idx=288 (list end) \\
27 & ...packy×5封禁/silicon余额不足 & idx=107 (list item) \\
26 & ...Codex 接入 packy、packy 密钥账本 & idx=163 (list item) \\
\bottomrule
\end{tabular}
\end{table}

Zero of the top-5 facts are about packy. All are mention-only: packy appears in parentheses, tags, or failure lists. The facts that \emph{are} about packy (e.g., ``packy分组机制: 一个key只能调本组模型'') have lower use\_count and are buried.

\subsection{Cause 2: Regex-Based Entity Extraction}

Table~\ref{tab:coverage} summarizes the extraction coverage. Entity extraction uses regex patterns for Chinese syntax (\texttt{X的Y}, \texttt{X支持Y}) and first-phrase heuristics. This misses 41.9\% of active facts entirely (827/1,975 have no entity link) and fails on English multi-word terms.

\begin{table}[h]
\centering
\caption{Entity extraction coverage by query term.}
\label{tab:coverage}
\begin{tabular}{lccc}
\toprule
Query term & Facts containing & With entity link & Coverage \\
\midrule
MEMORY.md & 42 & 38 & 90.5\% \\
fuseeval & 72 & 13 & 18.1\% \\
Memory Court & 43 & 2 & 4.7\% \\
temporal accuracy & 7 & 0 & 0.0\% \\
\bottomrule
\end{tabular}
\end{table}

The failure is structural: the regex extractor processes Chinese patterns only. ``Memory Court'' matches no pattern and extracts nothing. ``fuseeval v3-referee全量...'' extracts the entire first phrase as the entity name, making exact-match lookup fail.

\section{Method: Entity-Aware Hybrid Retrieval}

\subsection{Subject-Proximity Score}

We score each entity-linked fact by the position of the query entity in its content:

\begin{equation}
\text{ss}(c, e) = \begin{cases}
1.0 & \text{idx}(e, c) \leq 30 \quad \text{(subject)} \\
0.7 & \text{idx}(e, c) \leq 80 \quad \text{(prominent)} \\
0.5 & \text{count}(e, c) \geq 2 \quad \text{(repeated)} \\
0.3 & \text{idx}(e, c) \leq 200 \quad \text{(mid-text)} \\
0.1 & \text{otherwise} \quad \text{(passing reference)}
\end{cases}
\end{equation}

\noindent where $\text{idx}(e,c)$ is the character position of entity $e$ in content $c$.

The final score combines subject proximity (primary), Q-value (learned quality), and normalized use\_count (popularity):
$$\text{score} = 0.7 \cdot \text{ss} + 0.2 \cdot q + 0.1 \cdot \min(\text{uc}, 100)/100$$

\subsection{Hybrid Gating}

Entity-first retrieval is high-precision but low-recall: it fails completely when the entity index lacks the query term. We gate on coverage:

\begin{algorithmic}
\IF{entity index has $\geq 3$ facts with $\text{ss} \geq 0.7$}
    \STATE return entity-ranked results (subject-proximity order)
\ELSE
    \STATE return baseline hybrid search results
\ENDIF
\end{algorithmic}

This guarantees no regression: when entity extraction works, we get high precision; when it fails, we fall back to the baseline.

\section{Evaluation}

Table~\ref{tab:results} presents the main comparison. With only 10 queries, statistical significance cannot be reliably assessed; we report the raw counts and note that the improvement direction is consistent across 7 of 10 queries.

\begin{table}[h]
\centering
\caption{Precision@10: baseline vs.\ entity-aware retrieval (manually labeled ground truth, 100 query--result pairs across 10 queries).}
\label{tab:results}
\begin{tabular}{lcc}
\toprule
Method & Correct@10 & Precision@10 \\
\midrule
Baseline (hybrid search) & 28/100 & 28\% \\
Entity+SP (subject-proximity) & \textbf{73/100} & \textbf{73\%} \\
\midrule
\textbf{Improvement} & \textbf{+45pp} & \\
\bottomrule
\end{tabular}
\end{table}

\begin{finding}[Entity-Aware Hybrid Retrieval]
Entity-aware retrieval improves precision@10 from 28\% to 73\% (+45pp) with manually labeled ground truth. The improvement concentrates on queries where the entity index has correct subject-position facts: packy (1$\to$10), memory\_hub (2$\to$10), fuseeval (2$\to$10), Memory Court (1$\to$10). On queries where extraction fails, entity-aware retrieval correctly falls back to the baseline. The only regression is temporal accuracy (1$\to$0), where no entity link exists.
\end{finding}

\subsection{Analysis}

The improvement concentrates on queries where the entity index has correct subject-position facts: packy (1$\to$10), memory\_hub (2$\to$10), fuseeval (2$\to$10), Memory Court (1$\to$10). On queries where extraction fails, the method correctly falls back to baseline search.


The remaining 27\% errors are of two types: (1)~queries with no correct entity extraction (supersession, temporal accuracy, MEMORY.md); (2)~facts that mention the entity in subject position but are not the most current version.

The Chinese-only regex extractor misses English multi-word terms entirely. We add four deterministic English patterns: (1)~multi-word Title Case; (2)~acronyms; (3)~file names; (4)~lowercase technical terms followed by Chinese. Applied to the 827 facts with no entity link, these patterns add 4,136 new entity links. The enhanced index recovers additional entity links: Memory Court ($5\to$9) and supersession ($0\to$2) improve because their English terms now match the entity index. However, the overall precision@10 remains 73\% because gains on these queries are offset by the temporal accuracy regression ($1\to$0) and the net effect on other queries is neutral. The enhanced index improves entity coverage without degrading aggregate precision. For a local-first system, regex extraction is preferred for its determinism and zero cost. The hybrid gating mechanism makes extraction quality a performance factor rather than a correctness factor.

\paragraph{Generalizability.} The about-vs-mention problem affects any retrieval system that indexes text by surface features. Our subject-proximity score requires only the fact content and the query string---no language-specific parsing, no embedding model. It can be applied as a post-processing step on any existing retriever.

\paragraph{Limitations.} Our evaluation uses 10 queries from a single production deployment. The ground-truth labeling is based on manual inspection of each query--result pair, which is subjective but independent of the retrieval methods being compared. The entity index is built by a regex extractor with known coverage gaps; a system with better extraction would see larger gains from entity-first retrieval.

\section{Related Work}

Agent memory systems (mem0~\cite{mem0}, Zep~\cite{zep}, Letta~\cite{letta}) use embedding-based retrieval and report benchmark accuracy but do not measure production retrieval precision. MemoryAgentBench~\cite{mab} evaluates agent memory on task success metrics. ConflictBank~\cite{conflictbank} addresses contradiction detection, a different problem. To our knowledge, no prior work quantifies or addresses the about-vs-mention distinction in agent memory retrieval.

\section{Reproducibility}

All measurements are reproducible from the production SQLite database. The evaluation script and retrieval implementation are available at \url{https://github.com/lanbass869-cell/cac-artifact}.

\begin{verbatim}
# Baseline precision@10 (ground truth: manual labeling)
python entity_aware_retrieval.py --query "packy" --mode baseline
# Expected: 1/10 ABOUT (manual ground truth)

# Entity-aware hybrid precision@10
python entity_aware_retrieval.py --query "packy" --mode hybrid
# Expected: 10/10 ABOUT (manual ground truth)

# Entity extraction coverage
sqlite3 memory.db "SELECT COUNT(*) FROM facts f WHERE status='active'
  AND NOT EXISTS (SELECT 1 FROM fact_entities fe WHERE fe.fact_uid = f.uid)"
# Expected: 827 (41.9% of 1975 active)
\end{verbatim}

\section{Conclusion}

We identified, quantified, and partially solved a retrieval failure mode in production agent memory: systems cannot distinguish facts about a query entity from facts containing the query string. Entity-aware hybrid retrieval improves precision@10 from 28\% to 73\% with one known regression (temporal accuracy), using only character-position heuristics and existing entity indices.

\begin{thebibliography}{5}

\bibitem[Chhikara et~al.(2025)]{mem0}
Chhikara, P. et~al.
\newblock Mem0: Building Production-Ready AI Agents with Scalable Long-Term Memory.
\newblock \emph{arXiv:2504.19413}, 2025.

\bibitem[Rasmussen et~al.(2025)]{zep}
Rasmussen, R. et~al.
\newblock Zep: A Temporal Knowledge Graph Architecture for Agent Memory.
\newblock \emph{arXiv:2501.13956}, 2025.

\bibitem[Packer et~al.(2023)]{letta}
Packer, C. et~al.
\newblock MemGPT: Towards LLMs as Operating Systems.
\newblock \emph{arXiv:2310.08560}, 2023.

\bibitem[Hu et~al.(2025)]{mab}
Hu, Y. et~al.
\newblock Evaluating Memory in LLM Agents via Incremental Multi-Turn Interactions.
\newblock \emph{arXiv:2507.05257}, 2025.

\bibitem[Su et~al.(2024)]{conflictbank}
Su, Z. et~al.
\newblock ConflictBank: A Benchmark for Evaluating the Influence of Knowledge Conflicts in LLM.
\newblock \emph{arXiv:2408.12076}, 2024.

\end{thebibliography}

\end{document}
